% ---
% id: receiver/def:catalogue
% kind: def
% title: "Component catalogue"
% summary: >-
%   A component catalogue is a finite set of entry types, each fixing a
%   support shape, a parameter domain and a rule that yields a local element.
%   The catalogue submodule declares the entries and proves that each maps its
%   whole domain into local elements.
% prominence: standard
% status: accepted
% uses: [def:embed]
% mentions: [receiver/component-catalogue/thm:catalogue-admissibility, receiver/component-catalogue/def:catalogue-entries]
% assumes: []
% refs: []
% provenance: newly-derived
% lean:
% ---

Route (iii) declares a receiver as a graph of standard components rather than
as a matrix. This section shows when such a graph determines the
\tref{instrument matrix}: it does, under a condition that is stated rather
than assumed.

\begin{definition}\label{def:catalogue}
A \term[A finite set of entry types from which a route-(iii) receiver is assembled.]{component catalogue} $\Cat$ is a finite set of entry types. An
\term[A kind of optical component: a support shape, a parameter domain and a rule giving a local element for each parameter tuple.]{entry type} fixes a support shape (which modes of $\Modes$ it acts on), a
parameter domain, and a rule assigning to every parameter tuple in that domain
a \tref{local element} in the sense of Definition~\ref{def:embed}. The
catalogue of this contract is the one of Definition~\ref{def:catalogue-entries}
(Table~\ref{tab:catalogue}) in the component-catalogue submodule, and
Theorem~\ref{thm:catalogue-admissibility} there states that every entry type maps its whole
declared domain into local elements.
\end{definition}
